% anhang.tex: Anhang zum Pentagame-Orakel
% Eingebunden im Hauptdokument mit % anhang.tex: Anhang zum Pentagame-Orakel
% Eingebunden im Hauptdokument mit % anhang.tex: Anhang zum Pentagame-Orakel
% Eingebunden im Hauptdokument mit % anhang.tex: Anhang zum Pentagame-Orakel
% Eingebunden im Hauptdokument mit \include{anhang} (mit % davor abschalten).

\addchap{Anhang: Starke und schwache Siege}\markboth{Anhang}{\mbox{}}\label{kap:anhang}

\begin{einleitung}
Dieser Anhang verbindet das Orakel mit einer Beobachtung aus dem \textit{Pentagame Compendium} (Teil~IV, Kapitel~12). Er ist für Leser gedacht, die die Wendungen auch auf dem Brett weiterverfolgen möchten. Für die Befragung braucht man ihn nicht.
\end{einleitung}

\addsec*{Starke und schwache Siege}

\begin{einleitung}
Jede Wendung bildet ein stumpfes oder ein spitzes Dreieck (siehe \hyperref[kap:muster]{„Die zwölf Wendungen“}). Auf diese Form kommt es an, sobald man auch die Reihenfolge betrachtet, in der die Figuren herauskommen. Liegen die Ziele der ersten beiden Figuren nebeneinander (die erste Richtung ist ±1), bleibt für die dritte nur eine Möglichkeit, ein spitzes Dreieck zu bilden, aber zwei für ein stumpfes. Liegen sie übereck (±2), ist es umgekehrt. Der seltenere Abschluss ist der schwerere; das Compendium nennt ihn den stärkeren Sieg. Stark ist ein Sieg also, wenn auf zwei benachbarte Figuren ein spitzes Dreieck folgt oder auf zwei übereck liegende ein stumpfes.

Die erste Richtung und die Form des Dreiecks stecken vollständig in der Wendung. Ob ein Sieg stark oder schwach ist, hängt daher allein von der Wendung seines Tripels ab. Mit der Tabelle im Kapitel \hyperref[kap:muster]{„Die zwölf Wendungen“} ergibt sich: Stark sind \hyperlink{muster:II}{II, Die Prüfung}, \hyperlink{muster:VI}{VI, Der Rückhalt}, \hyperlink{muster:VIII}{VIII, Die Abkürzung} und \hyperlink{muster:X}{X, Die Läuterung}; die übrigen acht ergeben einen schwachen Sieg.

Alle vier starken wechseln zwischen Kreis und Stern: Auf eine Richtung von ±1 folgt eine von ±2, oder umgekehrt. Keine Wendung, die zweimal dieselbe Richtung geht, führt zu einem starken Sieg.

Das Compendium schlug vor, diese Unterscheidung bei einem Unentschieden (3:3) als Feinwertung zu nutzen: Ein starker Sieg schlägt dann einen schwachen. Im Spiel hat sich das nicht durchgesetzt. Für das Orakel zeigt es aber, dass die Wendungen nicht nur Deutungen sind, sondern auch auf dem Brett eine eigene Gestalt haben.
\end{einleitung} (mit % davor abschalten).

\addchap{Anhang: Starke und schwache Siege}\markboth{Anhang}{\mbox{}}\label{kap:anhang}

\begin{einleitung}
Dieser Anhang verbindet das Orakel mit einer Beobachtung aus dem \textit{Pentagame Compendium} (Teil~IV, Kapitel~12). Er ist für Leser gedacht, die die Wendungen auch auf dem Brett weiterverfolgen möchten. Für die Befragung braucht man ihn nicht.
\end{einleitung}

\addsec*{Starke und schwache Siege}

\begin{einleitung}
Jede Wendung bildet ein stumpfes oder ein spitzes Dreieck (siehe \hyperref[kap:muster]{„Die zwölf Wendungen“}). Auf diese Form kommt es an, sobald man auch die Reihenfolge betrachtet, in der die Figuren herauskommen. Liegen die Ziele der ersten beiden Figuren nebeneinander (die erste Richtung ist ±1), bleibt für die dritte nur eine Möglichkeit, ein spitzes Dreieck zu bilden, aber zwei für ein stumpfes. Liegen sie übereck (±2), ist es umgekehrt. Der seltenere Abschluss ist der schwerere; das Compendium nennt ihn den stärkeren Sieg. Stark ist ein Sieg also, wenn auf zwei benachbarte Figuren ein spitzes Dreieck folgt oder auf zwei übereck liegende ein stumpfes.

Die erste Richtung und die Form des Dreiecks stecken vollständig in der Wendung. Ob ein Sieg stark oder schwach ist, hängt daher allein von der Wendung seines Tripels ab. Mit der Tabelle im Kapitel \hyperref[kap:muster]{„Die zwölf Wendungen“} ergibt sich: Stark sind \hyperlink{muster:II}{II, Die Prüfung}, \hyperlink{muster:VI}{VI, Der Rückhalt}, \hyperlink{muster:VIII}{VIII, Die Abkürzung} und \hyperlink{muster:X}{X, Die Läuterung}; die übrigen acht ergeben einen schwachen Sieg.

Alle vier starken wechseln zwischen Kreis und Stern: Auf eine Richtung von ±1 folgt eine von ±2, oder umgekehrt. Keine Wendung, die zweimal dieselbe Richtung geht, führt zu einem starken Sieg.

Das Compendium schlug vor, diese Unterscheidung bei einem Unentschieden (3:3) als Feinwertung zu nutzen: Ein starker Sieg schlägt dann einen schwachen. Im Spiel hat sich das nicht durchgesetzt. Für das Orakel zeigt es aber, dass die Wendungen nicht nur Deutungen sind, sondern auch auf dem Brett eine eigene Gestalt haben.
\end{einleitung} (mit % davor abschalten).

\addchap{Anhang: Starke und schwache Siege}\markboth{Anhang}{\mbox{}}\label{kap:anhang}

\begin{einleitung}
Dieser Anhang verbindet das Orakel mit einer Beobachtung aus dem \textit{Pentagame Compendium} (Teil~IV, Kapitel~12). Er ist für Leser gedacht, die die Wendungen auch auf dem Brett weiterverfolgen möchten. Für die Befragung braucht man ihn nicht.
\end{einleitung}

\addsec*{Starke und schwache Siege}

\begin{einleitung}
Jede Wendung bildet ein stumpfes oder ein spitzes Dreieck (siehe \hyperref[kap:muster]{„Die zwölf Wendungen“}). Auf diese Form kommt es an, sobald man auch die Reihenfolge betrachtet, in der die Figuren herauskommen. Liegen die Ziele der ersten beiden Figuren nebeneinander (die erste Richtung ist ±1), bleibt für die dritte nur eine Möglichkeit, ein spitzes Dreieck zu bilden, aber zwei für ein stumpfes. Liegen sie übereck (±2), ist es umgekehrt. Der seltenere Abschluss ist der schwerere; das Compendium nennt ihn den stärkeren Sieg. Stark ist ein Sieg also, wenn auf zwei benachbarte Figuren ein spitzes Dreieck folgt oder auf zwei übereck liegende ein stumpfes.

Die erste Richtung und die Form des Dreiecks stecken vollständig in der Wendung. Ob ein Sieg stark oder schwach ist, hängt daher allein von der Wendung seines Tripels ab. Mit der Tabelle im Kapitel \hyperref[kap:muster]{„Die zwölf Wendungen“} ergibt sich: Stark sind \hyperlink{muster:II}{II, Die Prüfung}, \hyperlink{muster:VI}{VI, Der Rückhalt}, \hyperlink{muster:VIII}{VIII, Die Abkürzung} und \hyperlink{muster:X}{X, Die Läuterung}; die übrigen acht ergeben einen schwachen Sieg.

Alle vier starken wechseln zwischen Kreis und Stern: Auf eine Richtung von ±1 folgt eine von ±2, oder umgekehrt. Keine Wendung, die zweimal dieselbe Richtung geht, führt zu einem starken Sieg.

Das Compendium schlug vor, diese Unterscheidung bei einem Unentschieden (3:3) als Feinwertung zu nutzen: Ein starker Sieg schlägt dann einen schwachen. Im Spiel hat sich das nicht durchgesetzt. Für das Orakel zeigt es aber, dass die Wendungen nicht nur Deutungen sind, sondern auch auf dem Brett eine eigene Gestalt haben.
\end{einleitung} (mit % davor abschalten).

\addchap{Anhang: Starke und schwache Siege}\markboth{Anhang}{\mbox{}}\label{kap:anhang}

\begin{einleitung}
Dieser Anhang verbindet das Orakel mit einer Beobachtung aus dem \textit{Pentagame Compendium} (Teil~IV, Kapitel~12). Er ist für Leser gedacht, die die Wendungen auch auf dem Brett weiterverfolgen möchten. Für die Befragung braucht man ihn nicht.
\end{einleitung}

\addsec*{Starke und schwache Siege}

\begin{einleitung}
Jede Wendung bildet ein stumpfes oder ein spitzes Dreieck (siehe \hyperref[kap:muster]{„Die zwölf Wendungen“}). Auf diese Form kommt es an, sobald man auch die Reihenfolge betrachtet, in der die Figuren herauskommen. Liegen die Ziele der ersten beiden Figuren nebeneinander (die erste Richtung ist ±1), bleibt für die dritte nur eine Möglichkeit, ein spitzes Dreieck zu bilden, aber zwei für ein stumpfes. Liegen sie übereck (±2), ist es umgekehrt. Der seltenere Abschluss ist der schwerere; das Compendium nennt ihn den stärkeren Sieg. Stark ist ein Sieg also, wenn auf zwei benachbarte Figuren ein spitzes Dreieck folgt oder auf zwei übereck liegende ein stumpfes.

Die erste Richtung und die Form des Dreiecks stecken vollständig in der Wendung. Ob ein Sieg stark oder schwach ist, hängt daher allein von der Wendung seines Tripels ab. Mit der Tabelle im Kapitel \hyperref[kap:muster]{„Die zwölf Wendungen“} ergibt sich: Stark sind \hyperlink{muster:II}{II, Die Prüfung}, \hyperlink{muster:VI}{VI, Der Rückhalt}, \hyperlink{muster:VIII}{VIII, Die Abkürzung} und \hyperlink{muster:X}{X, Die Läuterung}; die übrigen acht ergeben einen schwachen Sieg.

Alle vier starken wechseln zwischen Kreis und Stern: Auf eine Richtung von ±1 folgt eine von ±2, oder umgekehrt. Keine Wendung, die zweimal dieselbe Richtung geht, führt zu einem starken Sieg.

Das Compendium schlug vor, diese Unterscheidung bei einem Unentschieden (3:3) als Feinwertung zu nutzen: Ein starker Sieg schlägt dann einen schwachen. Im Spiel hat sich das nicht durchgesetzt. Für das Orakel zeigt es aber, dass die Wendungen nicht nur Deutungen sind, sondern auch auf dem Brett eine eigene Gestalt haben.
\end{einleitung}